\documentclass[a4paper,12pt]{article}

\usepackage[utf8]{inputenc}
\usepackage[T1]{fontenc}
\usepackage[italian]{babel}
\usepackage{amsmath, amssymb}
\usepackage{graphicx}
\usepackage{hyperref}
\usepackage{geometry}
\usepackage{booktabs}
\usepackage{float}
\geometry{margin=2.5cm}

\geometry{margin=2.5cm}

\title{Report 3 -- Swap(k, TIMES)}

\begin{document}

\maketitle

\newpage

\section{Disclaimer}
Negli esempi ho cambiato il parametro \texttt{k} (quante volte
risolvo il problema QUBO) con il parametro \texttt{TIMES} (quante volte chiamo QALS o l'annealer 
nei casi in cui non si utilizza QALS).

\section{QALS}

\subsection{\texttt{ksp\_1.txt}}

\begin{table}[H]
\centering
\begin{tabular}{lc}
\toprule
Numero totale di run & 10 \\
Best Weight & 101 \\
Best Profit & 198 \\
Avg Weight GAP & -52 .0\\
Avg Profit GAP & 30.2 \\
Avg fQ & -719521.67 \\
fQ Min Fuound & -1380069.0 \\
n of fQ Min Found & 1 \\
Items & \texttt{[1, 0, 0, 0, 0, 0, 1, 1, 0, 0]} \\
\bottomrule
\end{tabular}
\caption{Risultati con $\lambda = \displaystyle \frac{\sum_{i=1}^{n} p_i}{3}$}
\end{table}

\begin{table}[H]
\centering
\begin{tabular}{lc}
\toprule
Numero totale di run & 10 \\
Best Weight & 101 \\
Best Profit & 198 \\
Avg Weight GAP & -98.8 \\
Avg Profit GAP & -20.4 \\
Avg fQ & -129.51 \\
fQ Min Fuound & -229.94.85 \\
n of fQ Min Found & 1 \\
Items & \texttt{[1, 1, 1, 0, 0, 1, 1, 1, 0, 1]} \\
\bottomrule
\end{tabular}
\caption{Risultati con $\lambda = \displaystyle \frac{\sum_{i=1}^{n} p_i}{650 \cdot C}$}
\end{table}

\subsection{\texttt{ksp\_2.txt}}

\begin{table}[H]
\centering
\begin{tabular}{lc}
\toprule
Numero totale di run & 10 \\
Best Weight & 978 \\
Best Profit & 7740 \\
Avg Weight GAP & -8567.7 \\
Avg Profit GAP & -4773.7 \\
Avg fQ & 742120534297.07 \\
fQ Min Fuound & -486215642301.0 \\
n of fQ Min Found & 1 \\
\bottomrule
\end{tabular}
\caption{Risultati con $\lambda = \displaystyle \frac{\sum_{i=1}^{n} p_i}{3}$}
\end{table}

\begin{table}[H]
\centering
\begin{tabular}{lc}
\toprule
Numero totale di run & 10 \\
Best Weight & 978 \\
Best Profit & 7740 \\
Avg Weight GAP & -8719.1 \\
Avg Profit GAP & -5643.4 \\
Avg fQ & 3562405.49 \\
fQ Min Fuound & 1727774.13 \\
n of fQ Min Found & 1 \\
\bottomrule
\end{tabular}
\caption{Risultati con $\lambda = \displaystyle \frac{\sum_{i=1}^{n} p_i}{650 \cdot C}$}
\end{table}

\section{NON QALS}

\subsection{\texttt{ksp\_1.txt}}

\begin{table}[H]
\centering
\begin{tabular}{lc}
\toprule
Numero totale di run & 10 \\
Best Weight & 101 \\
Best Profit & 198 \\
Avg Weight GAP & 0 \\
Avg Profit GAP & 0 \\
Avg fQ & -1455540.67 \\
fQ Min Fuound & -1455540.67 \\
n of fQ Min Found & 10 \\
Items & \texttt{[1, 1, 0, 0, 0, 0, 1, 0, 0, 0]} \\
\bottomrule
\end{tabular}
\caption{Risultati con $\lambda = \displaystyle \frac{\sum_{i=1}^{n} p_i}{3}$}
\end{table}

\begin{table}[H]
\centering
\begin{tabular}{lc}
\toprule
Numero totale di run & 10 \\
Best Weight & 101 \\
Best Profit & 198 \\
Avg Weight GAP & -76.0 \\
Avg Profit GAP & -104.0 \\
Avg fQ & -330.85 \\
fQ Min Fuound & -330.85 \\
n of fQ Min Found & 10 \\
Items & \texttt{[1, 1, 0, 0, 0, 1, 1, 1, 0, 0]} \\
\bottomrule
\end{tabular}
\caption{Risultati con $\lambda = \displaystyle \frac{\sum_{i=1}^{n} p_i}{650 \cdot C}$}
\end{table}

In entrambi i casi i risultati sempre uguali e uguali pure a quelli di Gurobi.

\subsection{\texttt{ksp\_2.txt}}

\begin{table}[H]
\centering
\begin{tabular}{lc}
\toprule
Numero totale di run & 10 \\
Best Weight & 978 \\
Best Profit & 7740 \\
Avg Weight GAP & -23.0 \\
Avg Profit GAP & 1861.8 \\
Avg fQ & -10200710058.53 \\
fQ Min Fuound & -10200710468.33 \\
n of fQ Min Found & 1 \\
\bottomrule
\end{tabular}
\caption{Risultati con $\lambda = \displaystyle \frac{\sum_{i=1}^{n} p_i}{3}$}
\end{table}

\begin{table}[H]
\centering
\begin{tabular}{lc}
\toprule
Numero totale di run & 1000 \\
Best Weight & 978 \\
Best Profit & 7740 \\
Avg Weight GAP & 0 \\
Avg Profit GAP & 0 \\
Avg fQ & -54748.31 \\
fQ Min Fuound & -54748.31 \\
n of fQ Min Found & 10 \\
\bottomrule
\end{tabular}
\caption{Risultati con $\lambda = \displaystyle \frac{\sum_{i=1}^{n} p_i}{650 \cdot C}$}
\end{table}

Anche in questo caso i risultati trovati da due ottimizzatori corrispondono.

\end{document}